\section{Аналогия с физикой твёрдого тела}

В классической физике твёрдого тела атомы расположены на решётке Браве,
элементарная ячейка задаётся зоной Вигнера--Зейтца, колебания описываются фононами,
спектр --- зонами Блоха.
Если под атомами лежит среда элементарных носителей масштаба~$\caStepSpace$,
тот же словарь повторяется на более глубоком уровне:
центры --- ячейки~$\hv$ масштаба~$\ell_P$,
частицы --- первичные вихри и моды поля~$z$,
а вещество --- дефекты и составные узлы на вакуумной решётке.
Фононы вещества образуют второй кристалл над кипящим первым.

